%%%%%%%%%%%%%%%%%%%%%%%%%%%%%%%%%%%%%%%%%%%%%%%%%%%%%
%
%           EE474 Intro to C Programming
   % combined into notes pack	%<stream, c>

\documentclass{article}	%<nh>

\usepackage{graphicx}	%<shn>
\usepackage{hyperref}	%<shn>
\usepackage{listings}	%<shn>

	%<*hn>
\oddsidemargin 0in
\evensidemargin 0in
\textwidth 6.5in
\textheight 8.7in
\topmargin 0.2in


	%<*shn>
\linespread{1.0}

\begin{document}


	%<*>

	%<*shnc>
\begin{center}
{\Large Structure and Compilation of C-Programs}
\end{center}
	%<*>


%%%%** Section 1 
\section{C Programs in Multiple Files}

A C program is typically written in multiple files.

Simpler editing (?)

Helps modularity and multiple authors.

    
File: \framebox{main.c}
\lstset{language=c}
\lstset{basicstyle=\ttfamily, commentstyle=\textit}
\begin{lstlisting}[frame=trbl]{}
main()
  {
   ... code  ...
  }

fcn1()
  {
  ... code ...
  }
\end{lstlisting}


File: \framebox{functions.c}

\lstset{language=c}
\lstset{basicstyle=\ttfamily, commentstyle=\textit}
\begin{lstlisting}[frame=trbl]{}
fnc2()
  {
   ... code  ...
  }

fcn3()
  {
  ... code ...
  }

... etc ...
\end{lstlisting}



Execution starts with first line of {\tt main() }.

All functions (except {\tt main()}) require a {\bf function prototype}
in any file in which they will be called or defined. 

All functions are visible from any file (but prototype must be supplied in the file).

%%%%** Section 2 
\section{Include Files (.h files)}
Text can be read in to your C-Source file just prior to compilation. 

File: \framebox{A.c}

\lstset{language=c}
\lstset{basicstyle=\ttfamily, commentstyle=\textit}
\begin{lstlisting}[frame=trbl]{}
#include X.h
main()
  {
   ... code  ...
  }
 
\end{lstlisting}
File: \framebox{X.h}

\lstset{language=c}
\lstset{basicstyle=\ttfamily, commentstyle=\textit}
\begin{lstlisting}[frame=trbl]{}
/*
 hello, this is just a comment
*/
\end{lstlisting}

output:

\lstset{language=c}
\lstset{basicstyle=\ttfamily, commentstyle=\textit}
\begin{lstlisting}[frame=trbl]{}
/*
 hello, this is just a comment
*/
main()
  {
   ... code  ...
  }
 
\end{lstlisting}



Reasons:
\begin{itemize}
	\item {\tt \#include} text which should be {\bf the same} in multiple ``.c" files. 
	\item keep and edit just 1 copy.
	\item less typing and ensures consistency.
\end{itemize}


    \includegraphics[width=4.5in]{cprog_figs/includes_multi_protos.png}
    
    A single {\tt .h} file introduces the exact same function prototypes
    and data declarations into multiple files without having to maintain 
    separate identical copies in multiple files. 
    

%%%%** Section 3 
\section{Other C pre-processor functions}

%%%%** Section 3.1
\subsection{{\tt\#define }}

like the search-replace function in your word processor.

ex:	%<!s>

{\tt \#define  TRUE  1}

c pre-processor will replace all occurrences of {\tt TRUE} with {\tt 1}.



%%%%** Section 3.2
\subsection{\tt\#ifdef / \#endif}

C pre-processor can turn on-off  blocks of code before compilation.

ex:	%<!s>

\lstset{language=c}
\lstset{basicstyle=\ttfamily, commentstyle=\textit}
\begin{lstlisting}[]{}
#define LINUX 1
//#define WINDOWS 1

...

#ifdef WINDOWS
     ... windows-specific code ...
#endif 

#ifdef LINUX
     ... linux-specific code ...
#endif 
\end{lstlisting}



%%%%** Section 3.3
\subsection{C pre-processor overview}

\includegraphics[width=5.5in]{cprog_figs/c-building.png}	%<chn>
    
	%<*>

The C pre-processor applies some convenience features to your code before 
compilation.  Specifically, the {\tt \#ifdef} macro can turn on and off C-code
until {\tt \#endif} is encountered.   You might use this to define a constant only 
if it has not already been defined by some earlier statements in a {\tt \#include}'d .h file.  If {\tt \#includes} were processed before the {\tt \#ifdef}, their statements
are taken into account when evaluating {\tt \#ifdef}.
Code from files indicated in {\tt \#include} statements (typically {\tt .h } files) is merged into your code.  Finally, all {\tt \#defines} are evaluated and the 
substitutions are made in your code.   Most likely, you will not need to worry about
these complexities unless you are working on very big packages of code.   For 
ECE474, you {\bf  probably} do not need {\tt \#ifdef}s at all. 

	%<*>
%%%%** Section 4 
\section{Building a C program}
For each ``.c" file:
\begin{enumerate}
	\item Run C pre-processor on the ``.c" file
	\item Compile output into ``object" file.
	\begin{itemize}
		\item Incomplete machine language program
		\item ``.o" file for each ``.c" file
	\end{itemize}

\end{enumerate}

{\bf Then: } ``Link" all ``.o" files, and library files,
	to create ``executable" binary file. 
%%%%** Section 5 
\section{Variable Scope}

\begin{itemize}
	\item Variable declared  inside a function is {\it local}.
	\item Variable declared outside a function is visible anywhere 
inside that ``.c" file.
	\item Variable declared outside a function can also be visible in 
other ``.c" files --- if {\tt extern} is used.
	\item Variables declared with prefix {\tt extern} must be declared
outside a function in {\it another } ``.c" file.
\end{itemize}


%%%%** Section 6 
\section{Function Scope}
\begin{itemize}

\item
Functions cannot be used without a function prototype.

\item
Each ``.c" file must have a function prototype for each function which is used
in that ``.c" file. 

\item
{\tt \#include } can help manage your function prototypes. 

\end{itemize}

\end{document}	%<shn>



